Put \(p=(q_S,q_0)\) and \(\widehat p=(\widehat q_S,\widehat q_0)\).  First, the soundness conclusion for the successful EXT-ID runs gives, cell by cell,
\[
F_{x,y}(q_S,q_0)=\psi_{x,y}(M)
\qquad ((x,y)\in\mathcal I)
\]
for every \(M\in\mathfrak F_{\epsilon}(\mathcal G,T_0;q_S,q_0)\).  The definition of the lexicographically ordered full-kernel collection therefore gives
\[
\mathbf F(q_S,q_0)=\boldsymbol\psi(M).
\]
This step uses the EXT-ID soundness theorem; it does not define the causal target to equal the observed-law functional.

We next derive the stochastic expansion.  The empirical residuals are in the declared multinomial tangent spaces because
\[
\mathbf 1^{\mathsf T}(e_o-q_S)=0,
\qquad
\mathbf 1^{\mathsf T}(e_o-q_0)=0.
\]
Consequently the source derivatives in the influence-function definition may be applied to them.  Linearity gives the exact identities
\[
\begin{aligned}
D_S\mathbf F[\widehat q_S-q_S]
&=\frac1{n_S}\sum_{i=1}^{n_S}
D_S\mathbf F[e_{O_i^S}-q_S]
=\frac1{n_S}\sum_{i=1}^{n_S}\phi_S(O_i^S),\\
D_0\mathbf F[\widehat q_0-q_0]
&=\frac1{n_0}\sum_{j=1}^{n_0}
D_0\mathbf F[e_{O_j^0}-q_0]
=\frac1{n_0}\sum_{j=1}^{n_0}\phi_0(O_j^0).
\end{aligned}
\]
For a probability vector \(q\) on a finite set and its empirical table \(\widehat q_N\), direct calculation yields
\[
\mathbb E_q\|\widehat q_N-q\|_2^2
=\frac1N\sum_o q(o)\{1-q(o)\}
=\frac{1-\sum_o q(o)^2}{N}\le \frac1N.
\]
The inequality
\[
\mathbf 1\{\|R\|_2>u\}\le \|R\|_2^2/u^2
\]
then shows that \(\widehat q_S-q_S=O_P(n_S^{-1/2})\) and \(\widehat q_0-q_0=O_P(n_0^{-1/2})\).  The sample-share assumption implies that \(n_S\), \(n_0\), and \(n\) have the same order.

By differentiability of the returned finite-dimensional map at \(p\),
\[
\mathbf F(\widehat p)-\mathbf F(p)
=D_S\mathbf F[\widehat q_S-q_S]
+D_0\mathbf F[\widehat q_0-q_0]+r_n,
\qquad
\frac{\|r_n\|_2}{\|\widehat q_S-q_S\|_2+\|\widehat q_0-q_0\|_2}
\xrightarrow{P}0.
\]
Every population denominator is strictly positive by the positive-denominator assumption.  Because there are finitely many continuous denominators, there is a neighborhood of \(p\) on which they remain positive.  The empirical tables enter this neighborhood with probability tending to one by the preceding bound.  Thus the zero fallback in the totalized plug-in definition is used with probability tending to zero and does not alter the displayed first-order expansion.  Hence \(r_n=o_P(n^{-1/2})\), and multiplication by \(\sqrt n\) gives the asserted asymptotic linear representation.

It remains to identify its limit.  The influence vectors are centered:
\[
\mathbb E_{q_S}\phi_S(O^S)
=D_S\mathbf F\!\left[\sum_o q_S(o)(e_o-q_S)\right]=0,
\qquad
\mathbb E_{q_0}\phi_0(O^0)=0.
\]
Form one triangular row from the mutually independent vectors
\[
R_{n,i}=\frac{\sqrt n}{n_S}\phi_S(O_i^S),
\quad 1\le i\le n_S,
\qquad
R_{n,n_S+j}=\frac{\sqrt n}{n_0}\phi_0(O_j^0),
\quad 1\le j\le n_0.
\]
The state spaces are finite, so both influence vectors are bounded.  Since \(\sqrt n/n_S=O(n^{-1/2})\) and \(\sqrt n/n_0=O(n^{-1/2})\), for every \(u>0\) all summands have norm at most \(u\) for sufficiently large \(n\); the Lindeberg sum is therefore eventually zero.  Independence of the two samples and the sample-share limit give
\[
\sum_{\ell=1}^{n}\operatorname{Cov}(R_{n,\ell})
=\frac n{n_S}\operatorname{Var}_{q_S}(\phi_S)
+\frac n{n_0}\operatorname{Var}_{q_0}(\phi_0)
\longrightarrow
\lambda^{-1}\operatorname{Var}_{q_S}(\phi_S)
+(1-\lambda)^{-1}\operatorname{Var}_{q_0}(\phi_0)
=\Sigma.
\]
The cited multivariate Lindeberg--Feller theorem now yields convergence of the linear term to \(N(0,\Sigma)\).  The cited asymptotic-equivalence theorem absorbs the \(o_P(1)\) remainder.  This proves both the displayed expansion and the claimed limit law.